% Source-derived chapter 10: The $p$-adic comparison theorem
% Original source: standard-conjectures-abelian-fourfolds
\chapter{The $p$-adic comparison theorem}
\label{standard-conjectures-abelian-fourfolds-chapter-10}
\source{standard-conjectures-abelian-fourfolds}

\begin{remark}[Source section]
\label{standard-conjectures-abelian-fourfolds-chapter-10-source}
\source{standard-conjectures-abelian-fourfolds}
\source{standard-conjectures-abelian-fourfolds}
This chapter reproduces the corresponding section of the retrieved
LaTeX source, including its definitions, statements, examples, and
proofs. It has not yet been translated into Lean.
\notready
\end{remark}

% The source section title was: The $p$-adic comparison theorem
We keep notation from Theorem \ref{mainthm} and Proposition \ref{padicreduction}.  
The aim of this and next sections is to compare the two $\Q_p$-quadratic spaces of rank two
\[ (V_{B,p},q_{B,p}) \coloneqq (V_B,q_B)\otimes_\Q \Q_p \hspace{0.5cm}  \textrm{and}\hspace{0.5cm}  (V_{Z,p},q_{Z,p}) \coloneqq (V_Z,q_Z)\otimes_\Q \Q_p\]
and deduce from this study Theorem \ref{mainthm}  (via Proposition \ref{padicreduction}). 

The core of the proof is in the next sections, we give here some preliminary results. We start by recalling the $p$-adic comparison theorem (Theorem \ref{tsuji}) and then we apply it to our geometric situation.  

\

Recall that $K$ denotes the $p$-adic field over which the motive $M$ from Theorem \ref{mainthm} is defined.



\begin{theorem}
\source{standard-conjectures-abelian-fourfolds}
\notready\label{javier} \cite{Fonfon}
\source{standard-conjectures-abelian-fourfolds}
There are two integral $\Q_p$-algebras \[B_\crys \subset B_\dR\] the first endowed with  actions of the Galois group $\Gal_K$ and of the absolute Frobenius $\varphi$ and the second endowed with a discrete valuation,  
\[\nu: B_\dR \longrightarrow \Z \cup \{+\infty\}\]
hence in particular endowed with a decreasing filtration \[\Fil^i B_\dR=\{x \in B_\dR, \nu(x)\geq i\}\] verifying 
\[\Fil^i \cdot \Fil^j \subset \Fil^{i+j}.\] 
The $\Q_p$-algebra $ B_\dR$ contains $\overline{\Q}_p$, an algebraic closure of $\Q_p$, and the intersection $ B_\crys \cap \overline{\Q}_p$ is the biggest non-ramified extension of $\Q_p$ inside $\overline{\Q}_p$.

Finally, the following equality holds \cite[Theorem 5.3.7]{Fontp}
 \begin{equation}\label{eq:fontaine}(B_\crys^{\varphi=\id} \cap \Fil^0B_\dR)=\Q_p.\end{equation}
\source{standard-conjectures-abelian-fourfolds}
\end{theorem}

\begin{theorem}[\cite{FontMess,Falt,CF}]
\source{standard-conjectures-abelian-fourfolds}
\notready\label{tsuji} 
\source{standard-conjectures-abelian-fourfolds}
There is an  equivalence of $\Q_p$-linear rigid categories 
\[D : \REP \longrightarrow \MOD\]
between the category $\REP$  of crystalline  $\Gal_K$-representations and the category $\MOD$ of admissible filtered $\varphi$-modules (see Convention (\ref{Hyodo-Kato})). This equivalence verifies the following properties:
\begin{enumerate}
\item\label{tsujican} There is a canonical identification 
\source{standard-conjectures-abelian-fourfolds}
\[ V \otimes B_\crys=  D(V)\otimes B_\crys \]
which is $\Gal_K$-equivariant and $\varphi$-equivariant. Moreover, the induced isomorphism
\[ V \otimes B_\dR=  D(V)\otimes B_\dR \]
respects the filtrations.

\item The functor $D$ is given by
\[V \mapsto D(V)= (V \otimes B_\crys)^{\Gal_K}\]
and its inverse is given by
\[ [W\otimes B_\crys]^{\varphi=\id} \cap  \Fil^0[W \otimes B_\dR] \mapsfrom W.\]

\item The equivalence $D$  is compatible with the two realization functors\footnote{This last fact is already present in the original works in $p$-adic Hodge theory although they are not written down in the motivic language. A reference is 
\cite[Theorem 1.2]{Niziol}, which can be immediately translated in our setting by  \cite[Proposition 4.2.5.1]{Andmot}.   This fact appears also in \cite[3.4.5]{Andmot} and it is implicitly  used in   \cite[7.4.1]{Andmot}. Nowadays much more general results are known, for instance this comparison of realization functors is available for mixed motives (i.e. motives of varieties which are not necessarely smooth and projective), see \cite[4.15]{NiziFred}.}
$R_p ((\cdot)_{\vert K}): \CHM(W) \rightarrow \REP$   and $R_\HK (\cdot): \CHM(W)\rightarrow \MOD$, namely we have
\[R_\HK (\cdot) = D \circ R_p ((\cdot)_{\vert K}).\]
\end{enumerate}
\end{theorem} 
\begin{corollary}
\source{standard-conjectures-abelian-fourfolds}
\notready\label{cortsuji}
\source{standard-conjectures-abelian-fourfolds}
There are canonical identifications (commuting with the extra structures):
\begin{equation} \label{eq:tsuji}
\source{standard-conjectures-abelian-fourfolds}
 (V_{B,p},q_{B,p}) =[(V_{Z,p},q_{Z,p})\otimes B_\crys]^{\varphi=\id} \cap  \Fil^0[(V_{Z,p},q_{Z,p})\otimes B_\dR].
\end{equation}
\begin{equation}\label{chissa}
\source{standard-conjectures-abelian-fourfolds}
 (V_{B,p},q_{B,p}) \otimes B_\crys=  (V_{Z,p},q_{Z,p}) \otimes B_\crys.
\end{equation}
\begin{equation}\label{chissa2}
\source{standard-conjectures-abelian-fourfolds}
 (V_{B,p},q_{B,p}) \otimes B_\dR=  (V_{Z,p},q_{Z,p}) \otimes B_\dR.
\end{equation}
Moreover, under these identifications,  the  equality of $\Q_p$-algebras  
\begin{equation}\label{stessaalgebra}
\source{standard-conjectures-abelian-fourfolds}
 \End_{\Gal_K}(V_{B,p}) = \End_{\varphi,\Fil^*}(V_{Z,p}\otimes \Frac(W(k)))
\end{equation}
holds.
\end{corollary}
\begin{proof}
 Proposition \ref{tortura}(\ref{giorno}) gives the identification
\[  V_{Z,p}\otimes \Frac(W(k)) =  R_\crys (M_{\vert k}).\]
and we have the identification $R_\crys (M_{\vert k})= R_\HK (M)$, see Convention (\ref{Hyodo-Kato}).
On the other hand, we have the equality
\[V_{B,p}  =  R_p (M_{\vert K})\]
because of the comparison theorem between singular and $p$-adic cohomology.
Finally, Theorem \ref{tsuji}(3) implies
\[ D(R_p (M_{\vert K}))   =   R_\HK (M).\]
Altogether we have the equality
\[D(V_{B,p}) = V_{Z,p}\otimes \Frac(W(k)).\]
Notice that these identifications are compatible with the quadratic forms as realization functors and the equivalence $D$ are tensor functors.
Hence Theorem \ref{tsuji}(2) gives  (\ref{eq:tsuji}) and Theorem \ref{tsuji}(1) gives  (\ref{chissa}) and   (\ref{chissa2}).
The identification (\ref{stessaalgebra}) follows from the fact that $D$ is an equivalence of $\Q_p$-linear categories.
\end{proof}

The following proposition settles Theorem \ref{mainthm} in the case where $V_B$ is of type $(0,0)$. It turns out that this case is easier than the others as the quadratic spaces $q_{B,p}$ and $q_{Z,p}$ are not only isomorphic (as predicted by Proposition  \ref{padicreduction}) but also equal (through the identification (\ref{chissa})).


\begin{prop}
\source{standard-conjectures-abelian-fourfolds}
\notready\label{ordinarycase}
\source{standard-conjectures-abelian-fourfolds}
Suppose that the Hodge structure $V_B$ is of type $(0,0)$, then $q_{B,p}$ and $q_{Z,p}$ are isomorphic, hence Theorem \ref{mainthm}  holds true in this case.
\end{prop}
\begin{proof}
First notice that   Frobenius acts trivially on $V_{Z,p}$ because the latter is spanned by  algebraic classes. Second, the hypothesis on the Hodge types gives  
$\Fil^0(V_{Z,p} \otimes K)=V_{Z,p} \otimes K$. Hence the relation (\ref{eq:tsuji}) implies 
\[(V_{B,p},q_{B,p}) =(V_{Z,p},q_{Z,p})\otimes (B_\crys^{\varphi=\id} \cap \Fil^0B_\dR).\]
Using (\ref{eq:fontaine}) we deduce the equality $q_{B,p} = q_{Z,p}$. Theorem \ref{mainthm}  then follows using  Proposition \ref{padicreduction}.
\end{proof}
\begin{remark}
\source{standard-conjectures-abelian-fourfolds}
\notready\label{boris}\begin{enumerate}
\source{standard-conjectures-abelian-fourfolds}
\item This proposition (together with the arguments of the previous sections) gives a full proof of Theorem \ref{thmscht} for ordinary abelian fourfolds. Indeed, in the ordinary case, all Galois invariant classes lift to $(0,0)$-classes as it can be shown with Shimura--Taniyama  formula  \cite[Lemma 5]{TateBour}. See Appendix \ref{geometricexamples} for related discussions.
\item\label{duccio} The hypothesis that $V_B$ is of type $(0,0)$ corresponds to the only case where one can hope that the algebraic classes in $V_Z$ might be lifted to characteristic zero, in which case Theorem \ref{mainthm} would follow from the Hodge--Riemann relations. We find an amusing coincidence that this conjecturally easier  case corresponds to an easier $p$-adic analysis.
\source{standard-conjectures-abelian-fourfolds}
\item The proof of the above proposition shows that, under the comparison isomorphisms of Corollary \ref{cortsuji}, one has the equality of $\Q_p$-vector spaces $V_{B,p} = V_{Z,p}$. The Hodge conjecture  predicts   that actually also the two $\Q$-structures $V_B$ and $V_Z$ should coincide as well. Is it possible to show the equality $V_{B} = V_{Z}$ without assuming the Hodge conjecture? We do not know. 

An analogous question can be formulated in the $\ell$-adic setting. Consider an ordinary abelian variety $A$  (of any dimension) together with its canonical lifting $\tilde{A}$. Fix an algebraic class on $A$. Does it corresponds to a Hodge class on $\tilde{A}$? This is a priori weaker than the Hodge conjecture: we do not ask that the algebraic cycle does lift to an algebraic cycle. Notice that if the answer to the question was affirmative then one would have a proof of the standard conjecture of Hodge type for $A$.
\end{enumerate}
\end{remark}

Thank to Proposition \ref{ordinarycase}  we are reduced to the case where $V_B$ is of type $(-i,i),(i,-i),$ for a positive integer $i$.
 It turns out that the only case which is interesting for the application to abelian fourfold is $i=1$, see Remark \ref{altreesotiche}(\ref{altreesotiche2}). We decided to work with a general $i>0$ to keep the possibility of applying  Theorem \ref{mainthm} to varieties other than abelian fourfolds, but the reader might find useful to think of the case $i=1$ in what follows.
 
 \begin{assumption}
\source{standard-conjectures-abelian-fourfolds}
\notready\label{ipos}
\source{standard-conjectures-abelian-fourfolds}
From now on we will suppose that the Hodge structure $V_B$ is not of type $(0,0)$. Equivalently there is a well-defined positive integer $i$ such that $\Fil^i(V_Z\otimes K)$ is a $K$-line and $\Fil^{i+1}(V_Z\otimes K)=0$.
\end{assumption}
 \begin{lemma}
\source{standard-conjectures-abelian-fourfolds}
\notready\label{isotrop}
\source{standard-conjectures-abelian-fourfolds}
 Under Assumption \ref{ipos} the line 
 $\Fil^i(V_Z\otimes K)$ is isotropic.
\end{lemma}
\begin{proof}
As the quadratic form $q$ is motivic its de Rham realization must respect the filtration, so $R_{\dR}(q)(\Fil^i(V_Z\otimes K)) \subset \Fil^{2i}R_{\dR}(\one).$ As $i$ is positive $\Fil^{2i}R_{\dR}(\one) =0.$
\end{proof}

\begin{prop}
\source{standard-conjectures-abelian-fourfolds}
\notready\label{propendo}
\source{standard-conjectures-abelian-fourfolds}
Under  Assumption \ref{ipos} the
$\Q_p$-algebra  $\End_{\Gal_K}(V_{B,p})$
is a field $F$ such that  $[F:\Q_p]=2$. Moreover, for all $v\in V_{B,p}$ and all $f\in F$,  we have the equality
\begin{equation} \label{eq:nonso}
\source{standard-conjectures-abelian-fourfolds}
 q_{B,p}(f \cdot v)=N_{F/\Q_p}(f) q_{B,p}(v).
\end{equation}
\end{prop}
\begin{proof}
Consider the ($\Q_p$-points of the) orthogonal  group $G=\SO(q_{B,p})$. We claim that the $\Q_p$-algebra
\[F\coloneqq \Q_p[G]\subset \End (V_{B,p})\]
satisfies all the properties of the statement. First, notice that this algebra is commutative and has dimension two by construction.  

\

As the quadratic form is induced by an algebraic cycle, the Galois group $\Gal_K$ must act on $V_{B,p}$ through $G$. Hence we have the inclusions  
\[\Q_p[\Gal_K]\subset  F \subset  \End (V_{B,p}).\] 
Let us show that $\Q_p[\Gal_K]$ is not of dimension one. If it were so, the algebra $\End_{\Gal_K}(V_{B,p})$ would be isomorphic to $M_{2\times 2}(\Q_p)$. By   (\ref{stessaalgebra}) the algebra $\End_{\varphi,\Fil^*}(V_{Z,p}\otimes \Frac(W(k)))$ would also be $M_{2\times 2}(\Q_p)$, which would imply that $V_{Z,p}\otimes \Frac(W(k))$ would be decomposed into the sum of two isomorphic filtered $\varphi$-modules. This is impossible as it would in particular imply that the filtration on $V_{Z,p}\otimes \Frac(W(k))$ would be a one step filtration, hence contradicting Assumption  \ref{ipos}.

As the $\Q_p$-algebra $\Q_p[\Gal_K]$ is not of dimension one, we deduce that the  inclusion $\Q_p[\Gal_K]\subset  F$ is an equality. The commutator of $F$ being $F$ itself, we also have
$F=\End_{\Gal_K}(V_{B,p})$. 

\

Let us show that the algebra $F$ is not isomorphic to $\Q_p \times \Q_p$. If it were so, arguing as before, we would have a decomposition of 
filtered $\varphi$-modules
$(V_{Z,p}\otimes \Frac(W(k))= W \oplus W'$ and each of the lines $W, W'$ would be isotropic. On the other hand  the line $\Fil^i(V_Z\otimes K)$ is also isotropic (Lemma \ref{isotrop}), hence, we would have  the equality 
\[W\otimes K =\Fil^i(V_Z\otimes K)\]
after possibly replacing $W$ with $W'$.

Now, as $W$ must be admissible, for any nonzero vector $w$ of $W$, the scalar $\alpha$ such that $\varphi(w)=\alpha w$ has $p$-adic valuation equal to $i$.
On the other hand, as  $V_{Z,p}$ is spanned by algebraic classes,  there is a nonzero vector of $W$ which is fixed by $\varphi$. As $i\neq 0$ by Assumption  \ref{ipos}, we deduce a contradiction.

\

As it is not isomorphic to $\Q_p \times \Q_p$ the algebra $F=\Q_p[\SO(q_{B,p})]$ must be a quadratic field extension of $\Q_p.$ The vector space $V_{B,p}$ can be identified with $F$ and the action of $F$ on it can be identified with the (left) multiplication. Under this identification we have $N_{F/\Q_p}(f)= \det(f \cdot)$ and the relation  (\ref{eq:nonso})   follows.
\end{proof}
\begin{remark}
\source{standard-conjectures-abelian-fourfolds}
\notready
When the motive $M$ from Theorem \ref{mainthm} is  an exotic motive (Section 7), the action of $F$ is induced by algebraic correspondences. The Hodge conjecture predicts that this should always be the case. Indeed, arguing as in the proof of Proposition \ref{propendo}, one can check that $E=\Q[\SO(\Q)]$ acts on $V_B$ respecting the Hodge decomposition and that $F=E \otimes \Q_p.$
\end{remark}
\begin{corollary}
\source{standard-conjectures-abelian-fourfolds}
\notready\label{nonso}
\source{standard-conjectures-abelian-fourfolds}
Keep Assumption  \ref{ipos} and let $F$ be the field of  Proposition \ref{propendo}. Then $F$ acts on $V_{Z,p}$ and the equality  (\ref{eq:tsuji}) is $F$-equivariant. Moreover,  for all $v\in V_{Z,p} $ and all $f\in F$  we have the equality
\begin{equation} \label{eq:nonso2}
\source{standard-conjectures-abelian-fourfolds}
q_{Z,p}(f \cdot v)=N_{F/\Q_p}(f) q_{Z,p}(v).
\end{equation}
\end{corollary}
\begin{proof}
If one replaces $V_{Z,p}$ by $V_{Z,p}\otimes \Frac(W(k))$ the statement is a combination of  Corollary \ref{cortsuji}   and Proposition \ref{propendo}. As $F$ commutes with the action of $\varphi$ and $V_{Z,p} \subset V_{Z,p}\otimes \Frac(W(k))$ is precisely the space of vectors which are fixed by $\varphi$, we deduce that $V_{Z,p}$ is stable by $F$ and the statement follows.
\end{proof}
\begin{corollary}
\source{standard-conjectures-abelian-fourfolds}
\notready\label{norms}
\source{standard-conjectures-abelian-fourfolds}
Keep Assumption  \ref{ipos} and let $F$ be the field of  Proposition \ref{propendo}. The following statements are equivalent:
\begin{enumerate}
\item The quadratic forms $q_{B,p}$ and $q_{Z,p}$ are isomorphic.
\item There exists a pair of nonzero vectors  $v_B\in V_{B,p}$ and $v_Z\in V_{Z,p}$ such that $q_{B,p}(v_B)$ and $q_{Z,p}(v_Z)$ are equal in $\Q_p^*/N_{F/\Q_p}(F^*).$
\item For any  pair of nonzero vectors  $v_B\in V_{B,p}$ and $v_Z\in V_{Z,p}$ we have that $q_{B,p}(v_B)$ and $q_{Z,p}(v_Z)$ are equal in $\Q_p^*/N_{F/\Q_p}(F^*).$
\end{enumerate}
\end{corollary}
\begin{proof}
This is a formal consequence of  formulas (\ref{eq:nonso}) and (\ref{eq:nonso2}).
\end{proof}
\begin{remark}
\source{standard-conjectures-abelian-fourfolds}
\notready\label{conventionF}
\source{standard-conjectures-abelian-fourfolds}
By the very construction of $F$, there are two actions of $F$ on  $V_{Z,p}\otimes F$ and this allows us to write the decomposition
\[V_{Z,p}\otimes F= V_{Z,+}  \oplus  V_{Z,-}\]
where $V_{Z,+}$ is the line where the two actions coincide and $V_{Z,-}$ is the line where the two actions are permuted by the non-trivial element of $\Gal(F/\Q_p)$. Using (\ref{eq:nonso2}), we see that $V_{Z,+}$ and $V_{Z,-}$ are also the two isotropic lines of the hyperbolic plane $V_{Z,p}\otimes F$. Notice that there is also an analogous decomposition
\[V_{B,p}\otimes F= V_{B,+} \oplus V_{B,-}\]
with analogous properties and that these decompositions are respected by Corollary \ref{cortsuji}. 

Finally, as $\Fil^i V_{Z,p}\otimes K$ is an isotropic line (Lemma \ref{isotrop}) it must coincide with $V_{Z,+}$ or  $V_{Z,-}$ (after extension of scalars to a field containing $F$ and $K$). We decide that the identification of $F$ and $K$ with a subfield of $\overline{\mathbb{Q}}_p$ is made to have the equality
\[\Fil^i V_{Z,p}\otimes  \overline{\mathbb{Q}}_p = V_{Z,+} \otimes  \overline{\mathbb{Q}}_p.\]  
\end{remark}
